\documentclass[11pt]{article}
\usepackage[margin=1in]{geometry}
\usepackage{xcolor}
\usepackage{tcolorbox}
\usepackage{hyperref}
\usepackage{enumitem}
\setlist{noitemsep,topsep=0pt}
\usepackage{booktabs}
\usepackage{array}
\usepackage{amsmath}

% Define OSU orange color
\definecolor{osaorange}{RGB}{220,115,38}

% Configure hyperref colors
\hypersetup{
    colorlinks=true,
    linkcolor=blue,
    urlcolor=blue,
    citecolor=blue
}

\title{}
\author{}
\date{}

\begin{document}

% Orange header box with black outline
\begin{tcolorbox}[colback=osaorange,colframe=black,boxrule=2pt,arc=0pt,left=6pt,right=6pt,top=10pt,bottom=10pt]
\centering
\textcolor{white}{\textbf{\Large Week 8 In-Class Worksheet}}\\
\textcolor{white}{\textbf{\large Contingency Tables and Chi-Square Tests}}
\end{tcolorbox}

\vspace{8pt}

\noindent\textbf{Name:} \rule{8cm}{0.4pt} \hfill \textbf{Date:} \rule{4cm}{0.4pt}\\[0.5ex]
\textbf{Course:} H524 Introduction to Biostatistics\\
\textbf{Topic:} Chi-Square Tests, Fisher's Exact Test, Effect Sizes

\vspace{8pt}

\section*{Instructions}
Work with a partner to solve the following problems. Show all your work including formulas, calculations, and interpretations. You may use R to verify answers, but demonstrate understanding by showing the steps.

\vspace{12pt}

\section{Problem 1: Fisher's Exact Test}

A small clinical trial tests a new treatment for a rare disease. Results:

\begin{center}
\begin{tabular}{l|cc|c}
& \textbf{Improved} & \textbf{Not Improved} & \textbf{Total} \\
\hline
\textbf{Treatment} & 7 & 3 & 10 \\
\textbf{Control} & 2 & 8 & 10 \\
\hline
\textbf{Total} & 9 & 11 & 20 \\
\end{tabular}
\end{center}

\vspace{6pt}

\textbf{Part A:} Calculate the expected counts for a chi-square test.

\vspace{0.3cm}

\begin{tabular}{l|cc}
& \textbf{Improved} & \textbf{Not Improved} \\
\hline
\textbf{Treatment} & & \\[1cm]
\textbf{Control} & & \\[1cm]
\end{tabular}

\vspace{1.5cm}

\textbf{Part B:} Are the chi-square test assumptions satisfied? Explain.

\vspace{0.3cm}

Rule: All expected counts should be $\geq 5$

\vspace{2.5cm}

\textbf{Part C:} Based on your answer to Part B, which test should be used?

\vspace{2.5cm}

\textbf{Part D:} Use R to perform Fisher's exact test and report the p-value and odds ratio.

\vspace{0.3cm}

\textbf{R code:}

\vspace{2cm}

P-value (two-sided) =

\vspace{1cm}

Odds Ratio =

\vspace{1cm}

95\% CI for OR =

\vspace{1cm}

\textbf{Part E:} State your conclusion about treatment effectiveness.

\vspace{3cm}

\newpage

\section{Problem 2: Chi-Square Test of Independence}

A cohort study follows 500 smokers and 500 non-smokers for 10 years to assess heart disease incidence.

\begin{center}
\begin{tabular}{l|cc|c}
& \textbf{Heart Disease} & \textbf{No Heart Disease} & \textbf{Total} \\
\hline
\textbf{Smoker} & 75 & 425 & 500 \\
\textbf{Non-smoker} & 35 & 465 & 500 \\
\hline
\textbf{Total} & 110 & 890 & 1000 \\
\end{tabular}
\end{center}

\vspace{6pt}

\textbf{Part A:} State the hypotheses for a chi-square test of independence.

\vspace{2.5cm}

\textbf{Part B:} Calculate the expected counts for each cell.

\vspace{0.3cm}

Formula: $E_{ij} = \frac{(\text{Row } i \text{ total}) \times (\text{Column } j \text{ total})}{\text{Grand total}}$

\vspace{0.3cm}

\begin{tabular}{l|cc}
& \textbf{Heart Disease} & \textbf{No Heart Disease} \\
\hline
\textbf{Smoker} & & \\[1cm]
\textbf{Non-smoker} & & \\[1cm]
\end{tabular}

\vspace{1.5cm}

\textbf{Part C:} Calculate the chi-square test statistic.

\vspace{0.3cm}

$$\chi^2 = \sum_{i,j} \frac{(O_{ij} - E_{ij})^2}{E_{ij}} = $$

\vspace{3.5cm}

\textbf{Part D:} Find the p-value and state your conclusion. ($df = (r-1)(c-1)$)

\vspace{0.3cm}

$df = $

\vspace{0.5cm}

P-value =

\vspace{1cm}

Conclusion:

\vspace{2cm}

\newpage

\section{Problem 3: Effect Sizes - Relative Risk and Odds Ratio (OPTIONAL - Week 9 Preview)}

\textbf{Note:} This problem previews next week's content on effect sizes. Try it for extra practice, but it will not be on this week's assignment.

\vspace{0.3cm}

A cohort study follows 500 individuals exposed to a chemical and 500 unexposed individuals for 5 years to assess disease incidence.

\begin{center}
\begin{tabular}{l|cc|c}
& \textbf{Disease} & \textbf{No Disease} & \textbf{Total} \\
\hline
\textbf{Exposed} & 80 & 420 & 500 \\
\textbf{Unexposed} & 40 & 460 & 500 \\
\hline
\textbf{Total} & 120 & 880 & 1000 \\
\end{tabular}
\end{center}

\vspace{6pt}

\textbf{Part A:} Calculate the relative risk (RR).

\vspace{0.3cm}

\textbf{Formula:} $RR = \frac{\text{Risk in exposed}}{\text{Risk in unexposed}} = \frac{a/(a+b)}{c/(c+d)}$

\vspace{0.3cm}

Risk in exposed group =

\vspace{1.5cm}

Risk in unexposed group =

\vspace{1.5cm}

Relative Risk (RR) =

\vspace{1.5cm}

\textbf{Part B:} Interpret the relative risk. What does this value tell us?

\vspace{2.5cm}

\textbf{Part C:} Calculate the odds ratio (OR).

\vspace{0.3cm}

\textbf{Formula:} $OR = \frac{a \times d}{b \times c}$ (cross-product ratio)

\vspace{0.3cm}

OR =

\vspace{2.5cm}

\textbf{Part D:} Compare RR and OR. Are they similar or different? Why?

\vspace{2.5cm}

\textbf{Part E:} Using R, verify your calculations and obtain a 95\% confidence interval for the odds ratio.

\vspace{0.3cm}

\textbf{R code:}

\vspace{2cm}

95\% CI for OR =

\vspace{1cm}

\textbf{Part F:} Based on the confidence interval, is there significant evidence of an association between exposure and disease? Explain.

\vspace{2.5cm}

\newpage

\section*{R Code Reference}

\textbf{Fisher's exact test:}
\begin{verbatim}
data <- matrix(c(7, 3, 2, 8), nrow=2, byrow=TRUE)
fisher.test(data)
\end{verbatim}

\textbf{Chi-square test of independence:}
\begin{verbatim}
data <- matrix(c(75, 425, 35, 465), nrow=2, byrow=TRUE)
chisq.test(data)
\end{verbatim}

\textbf{Effect sizes (Week 9 Preview):}
\begin{verbatim}
# Relative Risk
risk_exposed <- 80/500
risk_unexposed <- 40/500
RR <- risk_exposed / risk_unexposed

# Odds Ratio with CI
data <- matrix(c(80, 420, 40, 460), nrow=2, byrow=TRUE)
fisher.test(data)
\end{verbatim}

\end{document}
